\subsection{Non-recovery along canonical truncations}
The defect of Theorem~\ref{thm:defect} is defined relative to the known
intrinsic classes; a priori its elements could still be recovered along
every sequence.  For the canonical truncations this is not the case.

\begin{proposition}[Non-recovery]\label{prop:nonrecovery}
Let $T$, $I$ and $f$ be as in Theorem~\ref{thm:defect}, and let
$f_n:=\nabla p(x^{(n)})$, $x^{(n)}=z^{(n)}+Ue$, be the canonical truncations
of $f$.  Then:
\begin{enumerate}
\item[(a)] $f_n\in\NA\cap S_{p^*}$ and $f_n\to f$;
\item[(b)] for large $n$ the normer $x'_n:=x^{(n)}/p(x^{(n)})$ of $f_n$ is
C-tame, with $Q_{n,1}=\{2\}$, $Q_{n,m}=\emptyset$ for $m\ne1$, and no
degenerate peaks;
\item[(c)] $\Cset(f_n)\subseteq\R y^{(n)}$, where
$y^{(n)}:=\lambda_0u_{2,1}-(\Phi_1(2)^2w^{(n)}_1(2)/C^{(n)}_1)R_1^*w^{(n)}_1
\to\lambda_0u_{2,1}$;
\item[(d)] $\Li_n\Cset(f_n)\subseteq\R u_{2,1}$.  In particular, for
$\emptyset\ne K_1\ne K'$ and $\rho\in(0,1]$, $\rho g_{K_1}\notin
\Li_n\Cset(f_n)$: the mates $g_{K_1}$ are not recovered along the canonical
truncations.
\end{enumerate}
\end{proposition}

\begin{proof}
(a) $z^{(n)}\in c_{00}$, $\|z^{(n)}\|_\infty\le1$ and $z^{(n)}_1=1=\sgn a_1$,
so $\nabla q(x^{(n)})=a$ and $f_n\in\NA\cap S_{p^*}$ by
Proposition~\ref{prop:smooth}(c); $z^{(n)}\to z$ coordinatewise, and
Proposition~\ref{prop:approximants} gives $f_n\to f$.

(b) Let $c'_n:=q(x'_n)=1/p(x^{(n)})\to q_0>0$.  Since $|u_{k,m}(x^{(n)})|
\le q(x^{(n)})=1$, $\sigma'_{n,m}\le c'_nm2^{-m}$.  Exactly as in the proof
of Lemma~\ref{lem:firstrow}, using (P2) and (P3) at $x^{(n)}\in X_0$: $1$ is a
peak of every block, $\vartheta'_{n,m}\Phi_m(k)\le c'_n\pi_{k,m}$, and every
$(k,m)\ne(2,1)$ is a peak with margin $\mu'_{k,m}\ge c'_n\rho_{l(k,m)}$ if
$(k,m)$ is non-exceptional with $k\ge2$, and $\mu'_{k,m}\ge c'_n/4$
otherwise.  As $\rho_{l(k,m)}=(2\Phi_m(k)/c_{l(1,m)})^{1/2}\ge(2\Phi_m(k))^{1/2}$
and $\Phi_m(k+1)\le\Phi_m(k)/2$, for $s<c'_n/4$ we get
$\sum_{\mu'_{k,m}<s}\Phi_m(k)\le\sum_{\Phi_m(k)<s^2/(2c'^2_n)}\Phi_m(k)\le
s^2/c'^2_n=o(s)$: (MS) holds.  Next,
$u_{2,1}(x^{(n)})=u_{2,1}(\zh)-\sum_{j\in K',j>n}u_{2,1}(j)\to0$, while
$\vartheta'_{n,1}\Phi_1(2)\to\vartheta_1\Phi_1(2)>0$
(Proposition~\ref{prop:continuity}); so for large $n$,
$|u_{2,1}(x'_n)|<\vartheta'_{n,1}\Phi_1(2)$ and $(2,1)$ is a strict
non-peak.  Finally $a\in c_{00}$ and the contact set of $x'_n$ is
$K'\cap[1,n]$, which is finite.

(c) By Theorem~\ref{thm:tamefibre} at $\eta=x'_n$,
$\Cset(f_n)\subseteq S(f_n)=(\ell_1(\{1\})\cap(x'_n)^\perp)+\R y^{(n)}=\R
y^{(n)}$, since $x'_n(1)\ne0$.  By Proposition~\ref{prop:continuity},
$w^{(n)}_1(2)\to w_1(2)=0$ and $C^{(n)}_1\to C_1>0$, so $y^{(n)}\to
\lambda_0u_{2,1}$.

(d) If $\gamma_ny^{(n)}\to g$, then $(\gamma_n)$ is bounded because
$\|y^{(n)}\|\to\lambda_0\|u_{2,1}\|>0$; along a subsequence $\gamma_n\to\gamma$
and $g=\gamma\lambda_0u_{2,1}$.  The last claim follows from
Theorem~\ref{thm:defect}.
\end{proof}

Thus, for this $T$, the mates recovered along \emph{every} sequence are
contained in the certificate line $\R u_{2,1}$: universal recovery fails, and
a recovery of the switching mates must use norm-attaining approximants whose
certificate spans grow, for instance by putting base mass on the contact
windows (Section~\ref{sec:engineered}).  By Remark~\ref{rem:tameNA}, along
any sequence of norm-attaining approximants with C-tame normers, recovery is
a linear-algebra design problem: $\Li_n\Cset(f_n)\subseteq\Li_nS(f_n)$.

\subsection{Exact resonances}

\begin{proposition}[Exact resonances need infinite one-sided sets]\label{prop:exact}
Let $f\in S_{p^*}$, $\Omega\in V^*$ with $D:=L^*\Omega\ne0$, and suppose that
for a sequence $s_n\to0^+$,
\[
 q^*(a+s_nD)\le s(2s_n)\quad\text{and}\quad N_m(w_m-s_n\Omega_m)\le s(2s_n)
 \quad(m\in I).
\]
Then $\sum_{j\notin F}(|D_j|-z_jD_j)=0$: off $F$, $D$ is supported in the
contact set $K$ with $z_jD_j=|D_j|$.  Since $D\in Y\setminus\{0\}$ is not
finitely supported, $F\cup K$ is infinite.  In particular there are no such
transfers when $a\in c_{00}$ and $K$ is finite \textup{(}e.g.\ at every
norm-attaining $f$\textup{)}.
\end{proposition}

\begin{proof}
By Lemma~\ref{lem:base} (dropping nonnegative terms),
$q^*(a+sD)\ge1+sD(\zh)+s\sum_{j\notin F}(|D_j|-z_jD_j)$ for $s>0$, and
$N_m(w_m-s\Omega_m)\ge\ip{w_m-s\Omega_m}{\zeta_m}/\sigma_m=1-s\ip{\Omega_m}
{\zeta_m}/\sigma_m$.  With $s=s_n$, multiply the first inequality by $q_0$
and the $m$-th by $\sigma_m$ and add: since $q_0D(\zh)=D(\xi)=\ip{\Omega}
{L^{**}\xi}=\sum_m\ip{\Omega_m}{\zeta_m}$ and $q_0+\sum_m\sigma_m=1$,
\[
 q_0s_n\sum_{j\notin F}(|D_j|-z_jD_j)\le s(2s_n)-1\le2s_n^2 .
\]
Divide by $s_n$ and let $n\to\infty$.  Each term $|D_j|-z_jD_j$ is
nonnegative and vanishes only if $D_j=0$ or $|z_j|=1$ with $z_jD_j=|D_j|$.
\end{proof}

\begin{corollary}[Two-sided linear decompositions]\label{cor:twosided}
Let $g\in\Cset(f)$ and suppose there are $\tau>0$ and
$(B_\pm,\Omega_\pm)\in\ell_1\times V^*$ with $g=B_\pm+L^*\Omega_\pm$,
$\max(q^*(a+tB_+),\|w+t\Omega_+\|)\le s(t)$ for $0\le t\le\tau$ and
$\max(q^*(a+tB_-),\|w+t\Omega_-\|)\le s(t)$ for $-\tau\le t\le0$.  If
$F\cup K$ is finite, then $B_+=B_-$, $\Omega_+=\Omega_-$, and
$g\in\cl\Cert(f)$.
\end{corollary}

\begin{proof}
For $0<t\le\tau$ average the decompositions of $f+tg$ and $f-tg$:
$f=(a+\frac t2D)+L^*(w-\frac t2\Omega_D)$ with $D:=B_+-B_-$,
$\Omega_D:=\Omega_--\Omega_+$, $L^*\Omega_D=D$, and both component norms are
at most $s(t)$ by convexity.  If $D\ne0$, Proposition~\ref{prop:exact} with
$s_n=t_n/2$ contradicts the finiteness of $F\cup K$.  Hence $D=0$, and
$\Omega_+=\Omega_-$ by injectivity of $L^*$.  The common decomposition makes
$g$ locally split, and Theorem~\ref{thm:locsplit} applies.
\end{proof}

\begin{remark}[Resonance types; the sign rule]\label{rem:resonancetypes}
The transfer in the proof of Corollary~\ref{cor:twosided} exists for every
mate and every pair of admissible decompositions at scales $\pm t$; on the
contacts, for instance, the cheaply used parts of $B_+$ and of $B_-$ have
opposite signs, so they add up in $D=B_+-B_-$.  (A general ``matching at
scale'' statement for all one-sided resources is only available as a
sketch and is not used.)  A mate outside $\mathcal
K(f)$ therefore needs \emph{resonances}: cheap transfers with large
one-sided mass at arbitrarily small scales.  Proposition~\ref{prop:exact}
shows that exact (scale-free) resonances require infinitely many contacts
(or near-flip coordinates of $F$ when $a\notin c_{00}$);
Theorem~\ref{thm:defect} realizes one.  Approximate resonances run through
one-sided block carriers: weak peaks and near-threshold strict non-peaks
at depth comparable to the scale.  For these we only record an elementary
fact: if $u_{k,m}(\xi)\ne0$, then $\sgn w_m(k)=\sgn u_{k,m}(\xi)$ (on $P_m$
by Lemma~\ref{lem:threshold}, off $P_m$ because $w_m(k)=C_m\zeta_m(k)/
(\Phi_m(k)^2\sigma_m)$).  Consequently a one-sided usage of $(k,m)$ pushes in
the direction $-\sgn u_{k,m}(\xi)$, and two near-threshold carriers whose
vectors differ by $o(\Phi)$ in the $\xi$-direction (for instance Mart\'in's
cross-block near duplicates with small perturbations) serve the same side
of $t=0$ and cannot switch with each other.  This is a heuristic guide; no
classification of approximate resonances is proved here.  When the gaps of
the carriers are not summable, Theorem~\ref{thm:carriers} puts the
corresponding mates in $\cl\Cert(f)$; the summable case and the weak-peak
case are open (Remark~\ref{rem:summablegaps}).
\end{remark}

